\def\tituloconcepto{Generador de números aleatorios cuántico}
% Si el capitulo se compila solo (sin main.tex), esta cabecera abre el documento.
% Antes de % Si el capitulo se compila solo (sin main.tex), esta cabecera abre el documento.
% Antes de % Si el capitulo se compila solo (sin main.tex), esta cabecera abre el documento.
% Antes de \input{cabecera} cada capitulo define \tituloconcepto.
\ifdefined\enmain\else
  \documentclass[11pt,a4paper]{article}
  \input{preambulo}
  \begin{document}
  \title{\tituloconcepto}
  \author{\autor}
  \date{\today}
  \maketitle
\fi
 cada capitulo define \tituloconcepto.
\ifdefined\enmain\else
  \documentclass[11pt,a4paper]{article}
  % Preambulo compartido por todos los documentos. Compilar con XeLaTeX/LuaLaTeX
% (o Tectonic). Si se usa Overleaf, elegir el compilador XeLaTeX.
\usepackage[spanish,es-noshorthands]{babel}
\usepackage{amsmath,amssymb}
\usepackage{booktabs,array}
\usepackage[a4paper,margin=2.5cm]{geometry}
\usepackage{xcolor}
\usepackage{listings}
\usepackage[most]{tcolorbox}
\usepackage{quantikz}
\usepackage[hidelinks]{hyperref}

% ---- Autor (editar a gusto) -------------------------------------------------
\newcommand{\autor}{Genesis Silesky}

% ---- Notacion de Dirac y atajos ---------------------------------------------
% quantikz ya define \ket y \bra; DeclareRobustCommand los redefine sin error.
\DeclareRobustCommand{\ket}[1]{\left|#1\right\rangle}
\DeclareRobustCommand{\bra}[1]{\left\langle#1\right|}
\DeclareRobustCommand{\braket}[2]{\left\langle#1\middle|#2\right\rangle}
\newcommand{\C}{\mathbb{C}}
\newcommand{\R}{\mathbb{R}}
\DeclareMathOperator{\Tr}{Tr}

% ---- Cajas de apoyo --------------------------------------------------------
% "concepto": definiciones previas que hacen falta para entender el capitulo.
\newtcolorbox{concepto}[1]{breakable,colback=blue!4,colframe=blue!55!black,
  fonttitle=\bfseries,title={Concepto previo: #1}}
% "nota": aclaraciones y advertencias.
\newtcolorbox{nota}[1]{breakable,colback=orange!6,colframe=orange!70!black,
  fonttitle=\bfseries,title={#1}}

% ---- Estilos de listados ----------------------------------------------------
\lstdefinestyle{py}{language=Python,basicstyle=\ttfamily\footnotesize,
  breaklines=true,postbreak=\mbox{\textcolor{gray}{$\hookrightarrow$}\space},
  numbers=left,numberstyle=\tiny\color{gray},frame=single,
  commentstyle=\color{green!40!black}\itshape,
  keywordstyle=\color{blue!70!black}\bfseries,
  stringstyle=\color{red!60!black},showstringspaces=false,tabsize=4,
  columns=fullflexible,keepspaces=true}
\lstdefinestyle{salida}{basicstyle=\ttfamily\footnotesize,breaklines=true,
  frame=leftline,backgroundcolor=\color{gray!8},columns=fullflexible,
  keepspaces=true}

  \begin{document}
  \title{\tituloconcepto}
  \author{\autor}
  \date{\today}
  \maketitle
\fi
 cada capitulo define \tituloconcepto.
\ifdefined\enmain\else
  \documentclass[11pt,a4paper]{article}
  % Preambulo compartido por todos los documentos. Compilar con XeLaTeX/LuaLaTeX
% (o Tectonic). Si se usa Overleaf, elegir el compilador XeLaTeX.
\usepackage[spanish,es-noshorthands]{babel}
\usepackage{amsmath,amssymb}
\usepackage{booktabs,array}
\usepackage[a4paper,margin=2.5cm]{geometry}
\usepackage{xcolor}
\usepackage{listings}
\usepackage[most]{tcolorbox}
\usepackage{quantikz}
\usepackage[hidelinks]{hyperref}

% ---- Autor (editar a gusto) -------------------------------------------------
\newcommand{\autor}{Genesis Silesky}

% ---- Notacion de Dirac y atajos ---------------------------------------------
% quantikz ya define \ket y \bra; DeclareRobustCommand los redefine sin error.
\DeclareRobustCommand{\ket}[1]{\left|#1\right\rangle}
\DeclareRobustCommand{\bra}[1]{\left\langle#1\right|}
\DeclareRobustCommand{\braket}[2]{\left\langle#1\middle|#2\right\rangle}
\newcommand{\C}{\mathbb{C}}
\newcommand{\R}{\mathbb{R}}
\DeclareMathOperator{\Tr}{Tr}

% ---- Cajas de apoyo --------------------------------------------------------
% "concepto": definiciones previas que hacen falta para entender el capitulo.
\newtcolorbox{concepto}[1]{breakable,colback=blue!4,colframe=blue!55!black,
  fonttitle=\bfseries,title={Concepto previo: #1}}
% "nota": aclaraciones y advertencias.
\newtcolorbox{nota}[1]{breakable,colback=orange!6,colframe=orange!70!black,
  fonttitle=\bfseries,title={#1}}

% ---- Estilos de listados ----------------------------------------------------
\lstdefinestyle{py}{language=Python,basicstyle=\ttfamily\footnotesize,
  breaklines=true,postbreak=\mbox{\textcolor{gray}{$\hookrightarrow$}\space},
  numbers=left,numberstyle=\tiny\color{gray},frame=single,
  commentstyle=\color{green!40!black}\itshape,
  keywordstyle=\color{blue!70!black}\bfseries,
  stringstyle=\color{red!60!black},showstringspaces=false,tabsize=4,
  columns=fullflexible,keepspaces=true}
\lstdefinestyle{salida}{basicstyle=\ttfamily\footnotesize,breaklines=true,
  frame=leftline,backgroundcolor=\color{gray!8},columns=fullflexible,
  keepspaces=true}

  \begin{document}
  \title{\tituloconcepto}
  \author{\autor}
  \date{\today}
  \maketitle
\fi


\section{Generador de números aleatorios cuántico (QRNG)}

\subsection{Objetivo}
Generar bits aleatorios midiendo qubits en superposición y usarlos para
elegir la jugada de un bot de Tic-Tac-Toe, señalando las limitaciones de la
simulación.

\subsection{Conceptos previos}
\begin{concepto}{aleatoriedad, uniformidad y módulo}
\begin{itemize}
  \item \textbf{Aleatoriedad pseudoaleatoria vs.\ física:} un PRNG es un
    algoritmo determinista a partir de una semilla; una fuente física
    (p.\,ej.\ la medición de un qubit real) no lo es.
  \item \textbf{Uniformidad:} todos los resultados posibles tienen igual
    probabilidad.
  \item Medición y regla de Born (capítulo~04).
\end{itemize}
\end{concepto}

\subsection{Fundamento teórico}
$H\ket0=\tfrac1{\sqrt2}(\ket0+\ket1)$: medir da $0$ o $1$ con probabilidad
$\tfrac12$. Con $n$ qubits en paralelo se obtienen $n$ bits independientes
$\to$ un entero uniforme en $\{0,\dots,2^n-1\}$.

\paragraph{Sesgo por módulo.} Para elegir entre $m$ opciones con $x\bmod m$ la
distribución solo es uniforme si $m$ divide a $2^n$. En el código se eligen 4
bits ($16$ valores) para $m=5$ casillas; $16=3\cdot5+1$, así que la casilla
$0$ sale con probabilidad $4/16$ y las demás con $3/16$. Una solución es el
\emph{muestreo por rechazo}: descartar los valores $\ge 15$ y repetir.

\paragraph{La advertencia importante.} \texttt{AerSimulator} calcula las
probabilidades y luego toma muestras con un generador pseudoaleatorio. Por lo
tanto, \textbf{en simulación los bits no son aleatoriedad física}. Solo
ejecutando en hardware real la fuente es física (aun así, el ruido del equipo
puede introducir sesgo y se recomienda post-procesamiento/extractores).

\subsection{Código (completamente comentado)}
\lstinputlisting[style=py]{../09_qrng.py}

\subsection{Explicación celda por celda}
\begin{description}
  \item[Celda 2] La función arma un circuito con $n$ qubits, aplica $H$ a todos
    y mide. \texttt{shots=1, memory=True} devuelve la cadena de esa única
    ejecución.
  \item[Celda 3] Convierte los bits en un entero (\texttt{int(bits, 2)}) y
    elige una casilla libre del tablero con el operador módulo.
\end{description}

\subsection{Salida obtenida}
\lstinputlisting[style=salida]{../salidas/09_qrng.txt}
Distinta en cada ejecución (es el punto del generador).

\subsection{Interpretación}
Se ilustra el vínculo entre el concepto (medición) y una aplicación
concreta (un bot que elige jugadas). El ejemplo también expone una buena
práctica de ingeniería: validar que el mapeo de bits a opciones no introduzca
sesgo.

\subsection{Ejercicios}
\begin{enumerate}
  \item Implemente el muestreo por rechazo y verifique con 10\,000 jugadas que
    las 5 casillas tienen frecuencia $\approx20\%$.
  \item Mida con una prueba de $\chi^2$ la uniformidad de 10\,000 bits.
  \item Cambie el simulador por un backend real (sin ejecutar en la charla) y
    compare.
\end{enumerate}

\ifdefined\enmain\else
  \end{document}
\fi

